\section*{Hoofdstuk \ref{ch:b2:continuous-reactors} --- Continue reactoren en industriële processen}

\begin{solution}{exo:b2:continuous-reactors:1}
$\tau = V/Q_v = 2000/0.50 = \qty{4.0e3}{s}$, dus \qty{1.1}{h}.
\end{solution}

\begin{solution}{exo:b2:continuous-reactors:2}
$k\tau = 1.2 \times 10^{-3} \times 1800 = 2.16$. Tank: $X = 2.16/3.16 = 0.68$;
buis: $X = 1 - \mathrm e^{-2.16} = 0.88$.
\end{solution}

\begin{solution}{exo:b2:continuous-reactors:3}
Buis: $\tau = \ln 100/0.020 = \qty{230}{s}$. Tank: $X/(1 - X) = 99 = k\tau$,
$\tau = \qty{4950}{s}$, meer dan twintig keer langer: bij hoge \omterm{def:b2:continuous-reactors:conversion}{conversie} betaalt de
geroerde tank duur voor het werken bij de uitlaatconcentratie.
\end{solution}

\begin{solution}{exo:b2:continuous-reactors:4}
Per liter komt $120 \times 0.50 = \qty{60}{kJ}$ vrij, die
\qty{4.2}{kJ/K} opwarmt: $\Delta T_{\text{ad}} = \qty{14}{K}$. Een uitval van de koeling
zou het mengsel merkbaar maar niet gevaarlijk opwarmen, tenzij het oplosmiddel
dicht bij zijn kookpunt is of er in dat gebied een tweede reactie op gang komt.
\end{solution}

\begin{solution}{exo:b2:continuous-reactors:5}
Tank: $Da = X/(1 - X)^2 = 0.8/0.04 = 20$, $\tau = 20/0.010 = \qty{2000}{s}$.
Buis: $Da = X/(1 - X) = 4$, $\tau = \qty{400}{s}$. Verhouding 5, tegenover $4/\ln 5 =
2.5$ voor orde 1: hoe hoger de orde, hoe meer de geroerde tank lijdt onder
zijn lage concentratie.
\end{solution}

\begin{solution}{exo:b2:continuous-reactors:6}
$k\tau = 5.0$. Drie tanks: $X = 1 - (1 + 5/3)^{-3} = 0.947$. Eén tank: $5/6 =
0.833$. Buis: $1 - \mathrm e^{-5} = 0.993$.
\end{solution}

\begin{solution}{exo:b2:continuous-reactors:7}
$X = 0.70$ betekent $k\tau = 0.70/0.30 = 2.33$. Een verdubbeld debiet halveert $\tau$:
$k\tau = 1.17$, $X = 0.54$. Om weer 70\,\% te halen moet de \omterm{def:b2:continuous-reactors:residence-time}{verblijftijd}, en dus het
volume, verdubbeld worden.
\end{solution}

\begin{solution}{exo:b2:continuous-reactors:8}
Vrijgekomen: $80 \times 2.0 \times 1.0 \times 0.90 = \qty{144}{kW}$. Meegevoerd door de
uitlaatstroom, opgewarmd van \qty{330}{K} tot \qty{350}{K}: $4.0 \times 1.0 \times 20 =
\qty{80}{kW}$. De mantel moet het verschil afvoeren, \qty{64}{kW}.
\end{solution}

\begin{solution}{exo:b2:continuous-reactors:9}
Lineaire afmetingen $\times 10$: volume $\times 1000$, wandoppervlakte $\times 100$, oppervlakte
per volume-eenheid $\div 10$. De vrijgekomen warmte groeit met het volume, de
door de wand afgevoerde warmte met de oppervlakte: de grote reactor voert per liter
tien keer minder warmte af bij hetzelfde temperatuurverschil.
\end{solution}

\begin{solution}{exo:b2:continuous-reactors:10}
Buis: $\tau = \ln(0.05/0.10)/(0.05 - 0.10) = \qty{13.9}{min}$, $c_B/c_{A,0} =
\frac{0.10}{-0.05}(\mathrm e^{-1.386} - \mathrm e^{-0.693}) = 0.50$. Tank: $\tau =
1/\sqrt{0.005} = \qty{14.1}{min}$, $c_B/c_{A,0} = 1.414/(2.414 \times 1.707) =
0.34$. In de tank worden verse A en kant-en-klare B in hetzelfde volume gemengd: wat
B blijft altijd lang genoeg om tot C te worden omgezet, terwijl wat A nauwelijks
binnen is; in de buis heeft elke schijf dezelfde voorgeschiedenis.
\end{solution}

\begin{solution}{exo:b2:continuous-reactors:11}
Een hogere temperatuur van koelmiddel (en voeding) verschuift de afvoerlijn naar
rechts. Het koude snijpunt stijgt langzaam tot de lijn raakt aan
de onderste bocht van de S; daarna verdwijnt de koude toestand en springt de tank
naar de hete tak (ontsteking). Wordt de temperatuur weer verlaagd, dan blijft de tank
op de hete tak tot de lijn raakt aan de bovenste bocht, en pas
dan valt hij terug (doving), bij een lagere temperatuur dan de ontsteking: een
hysteresislus.
\end{solution}

\begin{solution}{exo:b2:continuous-reactors:12}
Batch: 75\,\% van \qty{2.0}{mol/L} blijft over, $1.5 \times 100 = \qty{150}{kJ/L}$,
$\Delta T = 150/4.0 = \qty{37.5}{K}$ boven de werktemperatuur, genoeg om
een kookpunt of een ontleding te bereiken. Semibatch: hoogstens
\qty{0.10}{mol/L} onomgezet, $\Delta T = 10/4.0 = \qty{2.5}{K}$. Langzaam toevoegen
begrenst de energie die in de reactor opgeslagen is.
\end{solution}

\begin{solution}{pb:b2:continuous-reactors:1}
\textbf{1.} $r = k[\text{ester}][\ce{OH-}]$; gelijke voedingen en een stoichiometrie
1 : 1 houden de twee concentraties gelijk, $r = kc^2$.
\textbf{2.} $1/c - 1/c_0 = kt$ met $c = 0.10c_0$: $kc_0t = 9$, $t = 9/(0.11
\times 0.10) = \qty{818}{s}$, ongeveer \qty{14}{min}.
\textbf{3.} $t_{1/2} = 1/(kc_0) = \qty{91}{s}$; bij orde 2 daalt de snelheid met
het kwadraat van de concentratie, zodat de tijd om die te halveren afhangt van het
vertrekpunt.
\textbf{4.} $Da = kc_0\tau = 0.011\,\tau$ ($\tau$ in seconden).
\textbf{5.} $Q_vc_0 - Q_vc - kc^2V = 0$, dus $c_0 - c = k\tau c^2$.
\textbf{6.} Met $c = c_0(1 - X)$: $c_0X = k\tau c_0^2(1 - X)^2$.
\textbf{7.} $Da = 0.90/0.10^2 = 90$.
\textbf{8.} $\tau = 90/0.011 = \qty{8.2e3}{s}$, \qty{2.3}{h}.
\textbf{9.} $V = 0.50 \times 8182 = \qty{4.1e3}{L}$, \qty{4.1}{m^3}.
\textbf{10.} De uitlaatconcentratie, $0.010\,\unit{mol/L}$: de snelheid is
overal honderd keer lager dan aan de inlaat.
\textbf{11.} $Da = X/(1 - X) = 9$, $\tau = \qty{818}{s}$, $V = \qty{409}{L}$.
\textbf{12.} Tien keer minder volume voor de buis.
\textbf{13.} $c_{j-1} - c_j = k\tau_1c_j^2$, met $\tau_1$ de \omterm{def:b2:continuous-reactors:residence-time}{verblijftijd} van
één tank.
\textbf{14.} $\tau_1 = 850/(3 \times 0.50) = \qty{567}{s}$, $k\tau_1 =
\qty{62.4}{L/mol}$. Oplossen van $62.4c_j^2 + c_j - c_{j-1} = 0$: $c_1 = 0.0328$,
$c_2 = 0.0163$, $c_3 = \qty{0.0100}{mol/L}$: 90\,\% \omterm{def:b2:continuous-reactors:conversion}{conversie}.
\textbf{15.} De buis, of de cascade (tien en vijf keer kleiner dan één
tank). Eén tank kan toch winnen door zijn uniforme temperatuur, eenvoudige
reiniging en regeling, en omdat volume hier goedkoop is.
\textbf{16.} $\Delta T_{\text{ad}} = 55\,000 \times 100/(997 \times 4180) =
\qty{1.3}{K}$ (\qty{1.2}{K} bij 90\,\%).
\textbf{17.} $55 \times 0.50 \times 0.10 \times 0.90 = \qty{2.5}{kW}$.
\textbf{18.} Nee: de uitlaatstroom voert de warmte af met een stijging van ongeveer
een kelvin; een runaway is onmogelijk.
\textbf{19.} Twintig keer geconcentreerder: de adiabatische stijging wordt
\qty{26}{K} en de warmte \qty{50}{kW}, wat een koelspiraal rechtvaardigt; omdat de snelheid $kc_0$
twintig keer hoger is, zou de tank voor 90\,\% twintig keer kleiner zijn
(\qty{0.20}{m^3}).
\textbf{20.} $0.50 \times 0.10 \times 0.90 = \qty{0.045}{mol/s}$, $\times 86\,400
\times 82.03 = \qty{319}{kg}$ natriumethanoaat per dag.
\textbf{21.} De snelheidsconstante groeit met de temperatuur (Arrhenius), dus het
volume krimpt; de prijs is het opwarmen van de voeding en, voor andere
reacties, een lagere \omterm{def:b2:continuous-reactors:selectivity}{selectiviteit} en meer damp.
\textbf{22.} De enkele geroerde tank moet $\boldsymbol{V \approx
\qty{4.1}{m^3}}$ bevatten.
\end{solution}
